% 题干转写：完成
% 解析转写：完成
% 图形核对：完成
\section{第40届复赛第3题：楔块上匀质圆球的滚动与滑动}

\subsection{题干}

如图\ref{fig:cpho40q03}，倾角为 $\theta$、质量为 $M$ 的三角形楔块放在光滑的水平地面上。其斜面上有一个质量为 $m$、半径为 $r$ 的匀质圆球；圆球从静止开始自由向下运动，整个过程中楔块不发生转动。重力加速度大小为 $g$。设圆球与楔块间的静摩擦因数和滑动摩擦因数均为 $\mu$。

\begin{figure}[htbp]
  \centering
  \begin{tikzpicture}[>=Stealth,scale=0.9]
  \draw[thick] (-4.0,-1.5) -- (4.0,-1.5);
  \fill[gray!15] (-2.9,-1.5) -- (2.7,-1.5) -- (2.7,1.18) -- cycle;
  \draw[thick] (-2.9,-1.5) -- (2.7,-1.5) -- (2.7,1.18) -- cycle;
  \draw[thick] (0.76,0.15) circle (0.57);
  \fill (0.76,0.15) circle (1.3pt) node[below left] {$C$};
  \draw[densely dashed] (0.76,0.15) -- (0.25,-0.34);
  \draw (2.15,-1.5) arc[start angle=180,end angle=137,radius=0.58];
  \node at (1.73,-1.10) {$\theta$};
  \node at (-1.45,-0.86) {楔块 $M$};
  \node[above] at (0.78,0.75) {圆球 $m,r$};
\end{tikzpicture}

  \caption{第40届复赛第3题图3a的TikZ重绘。}
  \label{fig:cpho40q03}
\end{figure}

\begin{enumerate}
  \item 若 $\mu$ 足够大，使圆球无滑动地滚下，求：
  \begin{enumerate}
    \item 楔块相对于地面的加速度大小 $a_0$；
    \item 圆球质心相对于楔块的加速度大小 $a_c$；
    \item 地面对楔块的支持力大小 $N$；
    \item 楔块对圆球的支持力大小 $N_1$；
    \item 使匀质圆球保持无滑动滚动所需的最小摩擦因数 $\mu_0$。
  \end{enumerate}
  \item 若 $\mu<\mu_0$，且楔块斜面足够长，求圆球由静止运动 $\Delta t$ 后，其与楔块接触点 $P$ 相对于楔块的速度。
\end{enumerate}

\subsection{参考解答整理}

取沿斜面向下为正方向。在相对于楔块静止的参考系中，圆球质心的运动方程、转动方程和无滑动约束分别为
\[
ma_c=ma_0\cos\theta+mg\sin\theta-f,
\qquad I_C\ddot\varphi=fr,
\]
\[
v_c=r\dot\varphi,
\qquad a_c=r\ddot\varphi,
\qquad I_C=\frac25mr^2.
\]
因此 $f=\frac25ma_c$。又由地面参考系中体系水平方向动量守恒，
\[
m(v_c\cos\theta-v_0)=Mv_0,
\qquad m(a_c\cos\theta-a_0)=Ma_0.
\]
\paragraph{（1）（i）楔块相对于地面的加速度。}
联立可得
\[
a_0=\boxed{\frac{5mg\sin\theta\cos\theta}
{7(M+m)-5m\cos^2\theta}}.
\]

\paragraph{（1）（ii）圆球质心相对于楔块的加速度。}
\[
a_c=\frac{5(M+m)g\sin\theta}
{7(M+m)-5m\cos^2\theta}.
\]

\paragraph{（1）（iii）地面对楔块的支持力。}
将楔块和圆球视为整体，竖直方向质心加速度为
\[
a_y=-\frac{ma_c\sin\theta}{M+m}.
\]
由竖直方向牛顿第二定律，以及垂直斜面方向的平衡关系，得到
\[
N=(M+m)g\,
\frac{M+\frac27m}{M+m-\frac57m\cos^2\theta}.
\]

\paragraph{（1）（iv）楔块对圆球的支持力。}
\[
N_1=mg\cos\theta\,
\frac{M+\frac27m}{M+m-\frac57m\cos^2\theta}.
\]
\paragraph{（1）（v）保持无滑动滚动的最小摩擦因数。}
从 $\mu_0=f/N_1$ 得到维持纯滚动的临界摩擦因数
\[
\mu_0=\frac{2(M+m)}{7M+2m}\tan\theta.
\]

\paragraph{（2）发生滑动后接触点的相对速度。}
当 $\mu<\mu_0$ 时，圆球发生滚动和滑动。此时 $f=\mu N_1$，联立圆球与楔块的关系，可得
\[
a_c=\frac{(M+m)g(\sin\theta-\mu\cos\theta)}
{M+m-m\cos^2\theta-\mu m\sin\theta\cos\theta}.
\]
接触点相对于楔块的速度为 $v_P=v_c-r\dot\varphi$，于是
\[
v_P=
\frac{(M+m)g\sin\theta-\mu\left(m+\frac72M\right)g\cos\theta}
{M+m-m\cos^2\theta-\mu m\sin\theta\cos\theta}\,\Delta t.
\]
当 $\mu=\mu_0$ 时上式给出 $v_P=0$，与纯滚动条件一致；当 $\mu<\mu_0$ 时，接触点相对楔块沿斜面向下运动。
